\section{Introduction}
\label{sec:introduction}

The geographic sorting of American partisan voters is the central confound in
partisan gerrymandering litigation. Democratic voters are concentrated in dense
urban cores---central cities, university towns, and inner suburbs---while Republican
voters are distributed more evenly across suburban, exurban, and rural areas.
This sorting is a consequence of residential choice, economic geography, and
social network dynamics, not mapmaker design. It would appear in any redistricting
plan, however drawn, because it reflects where people live, not how boundaries
are drawn.

The Mean-Median Difference captures this sorting pattern and provides a natural
tool for measuring how much of a map's partisan skew is attributable to geography
versus deliberate manipulation. As introduced by \citet{wang2016}, MM is:
\[
\text{MM} = \bar{v}_D - \tilde{v}_D,
\]
where $\bar{v}_D$ is the mean Democratic vote share across all $k$ districts
and $\tilde{v}_D$ is the median. When Democrats are concentrated in a few
high-margin districts, the mean is pulled above the median by those outlier
districts, producing positive MM. When the distribution is approximately symmetric,
mean and median are equal and MM is near zero.

The confound problem is direct: any positive MM might reflect either (1) the
natural geographic sorting of Democratic voters into a few dense urban districts,
or (2) deliberate packing of Democratic voters into fewer, higher-margin districts
by mapmakers. In federal litigation, opponents of partisan gerrymandering claims
routinely argued that positive MM is merely evidence of where Democrats choose to
live, not evidence of manipulation. Federal courts never resolved this dispute;
\textit{Rucho} made the resolution unnecessary for federal purposes.

State courts, however, must resolve it. A state constitutional claim that a map
``prioritizes partisan advantage'' (NC) or deviates from what ``neutral criteria
would produce'' (PA) requires a baseline against which enacted map properties
can be compared. \textsc{Bisect} algorithmic maps provide that baseline: they
are drawn from the same geographic constraints as enacted maps (the same census
tracts, the same adjacency structure, the same population distribution) but
without any partisan data. The MM of a \textsc{bisect} map therefore measures
the geographic sorting component---the baseline MM that the residential distribution
of voters produces regardless of mapmaker choices. The gap between enacted MM
and \textsc{bisect} MM measures the mapmaker contribution.

This paper develops that empirical argument. Section~\ref{sec:definition} defines MM
mathematically and documents its advantages over EG. Section~\ref{sec:behavior}
explains why \textsc{bisect} reduces MM and why the residual is attributable to
geography. Section~\ref{sec:empirical} presents multi-state empirical results with
election-cycle sensitivity. Section~\ref{sec:legal} documents MM's legal status
through 2026. Section~\ref{sec:conclusion} provides expert witness guidance.
